\section{Problem Definition and Observation--Action Safety Gap}
\label{sec:ieee-oasg}

ORIUS starts from a narrow safety object with broad consequences. Let $\zt$ denote the
physically true state, $\ztt$ the state exposed to the nominal controller, and
$\mathcal{A}(\cdot)$ the set of admissible actions under the relevant safety envelope.
The observation--action safety gap is the event
\[
\OASG_t = \mathbf{1}\{a_t \in \mathcal{A}(\ztt)\ \land\ a_t \notin \mathcal{A}(\zt)\},
\]
which means that the controller is internally consistent on the observed state while
the plant is externally unsafe on the true state. This is the failure mode that
reappears across battery dispatch, time-to-collision control, process envelopes,
clinical intervention logic, navigation corridors, and flight-energy bounds.

The defining difficulty is that OASG is not merely a state-estimation problem. A better
estimator helps, but the safety event is ultimately about actuation semantics. If the
runtime continues to treat the observed state as authoritative after freshness loss,
dropout, delayed updates, sensor disagreement, or reliability collapse, then even a
nominally optimal controller may produce a physically unsafe action. ORIUS therefore
places the universal object at the boundary between degraded observation and released
actuation rather than at the forecaster, planner, or dynamics model alone.

This framing also explains why a single universal policy is not the goal. Different
domains retain different state spaces, action geometries, intervention semantics, and
fallback policies. What can be shared is the runtime grammar: observation reliability,
inflated state uncertainty, tightened admissible actions, repaired actuation, and
certificate semantics. The paper argues that this runtime grammar is the right unit of
cross-domain reuse for Physical AI safety.

\begin{table}[t]
\centering
\caption{Canonical OASG runtime objects used throughout the IEEE draft.}
\label{tab:ieee-oasg-objects}
\small
\begin{tabularx}{\columnwidth}{@{}p{0.20\columnwidth}X@{}}
\toprule
\textbf{Object} & \textbf{Role in the runtime safety layer}\\
\midrule
$\ztt$ & Observed state or belief state exposed to the nominal controller.\\
$\zt$ & Physically true state whose envelope determines actual safety.\\
$w_t$ & Reliability assessment derived from freshness, missingness, and consistency.\\
$U_t$ & Observation-consistent uncertainty set widened when reliability drops.\\
$\mathcal{A}_t^{\mathrm{safe}}$ & Tightened admissible action set under degraded observation.\\
$a_t^{\mathrm{safe}}$ & Repaired or fallback action released by the runtime.\\
Certificate & Auditable payload recording reliability, margin, repair mode, and lifecycle state.\\
\bottomrule
\end{tabularx}
\end{table}

Two practical consequences follow. First, a runtime layer must expose the quality of
the observation surface as a first-class object rather than burying it inside model
residuals. Second, a universal benchmark must score safety on the true state rather
than on nominal controller legality alone. Those consequences drive the ORIUS runtime
design, the benchmark schema, and the claim boundary used throughout the rest of the
paper.
